\chapter{Introduction}
\begin{comment}
\section{Contexte et motivation}

La robotique mobile autonome nécessite une perception rapide et fiable de l’environnement pour naviguer en toute sécurité. Les caméras classiques acquièrent des images à une fréquence fixe, ce qui peut poser plusieurs problèmes :
\begin{itemize}
    \item Latence élevée dans des environnements dynamiques.  
    \item Charge de calcul importante pour le traitement d’images haute fréquence.  
    \item Difficulté à gérer des scènes avec une grande variation de luminosité (plage dynamique limitée).  
\end{itemize}

Les capteurs neuromorphiques (ou caméras à événements), inspirés de la rétine humaine, adoptent un principe différent : chaque pixel détecte les changements d’intensité de manière indépendante et émet un événement uniquement lorsqu’un changement significatif se produit.  

\textbf{Avantages de cette technologie :}
\begin{itemize}
    \item \textbf{Latence extrêmement faible :} les événements sont générés en temps réel, offrant une perception quasi instantanée.  
    \item \textbf{Faible consommation énergétique :} seulement les pixels qui détectent un changement émettent des événements, ce qui réduit la charge globale.  

    \item \textbf{Plage dynamique élevée :} ces caméras peuvent fonctionner dans des conditions très contrastées, de l’ombre à la lumière intense.  
    \item \textbf{Robustesse aux mouvements rapides :} les flux d’événements permettent de suivre des objets en mouvement rapide sans flou de mouvement.  
\end{itemize}

Ces caractéristiques rendent les caméras neuromorphiques particulièrement adaptées aux robots mobiles embarqués, qui doivent réagir rapidement aux changements dans l’environnement tout en limitant la consommation d’énergie.  

\section{Problématique}

Malgré leurs avantages, les capteurs neuromorphiques posent des défis spécifiques :
\begin{itemize}
    \item Les données produites sont des flux d’événements asynchrones, ce qui nécessite de repenser les algorithmes classiques de traitement d’image.  
    \item La conversion en images traditionnelles peut entraîner une perte d’information temporelle.  
    \item Les algorithmes de détection d’obstacles et d’égomotion doivent être adaptés pour exploiter la nature spatio-temporelle des événements.  
    \item La calibration et l’intégration sur une plateforme robotique mobile nécessitent une configuration spécifique.  
\end{itemize}

Ces contraintes motivent le développement de méthodes de perception spécifiques aux caméras DVS pour la navigation autonome.  

\section{Objectifs du projet}

L'objectif principal de ce projet est de développer et d’évaluer des algorithmes de perception adaptés aux caméras neuromorphiques pour la navigation autonome d’un robot mobile (robot Limo).  

\subsection{Objectifs spécifiques}

\begin{itemize}
    \item Développer des méthodes de détection d’obstacles exploitant les flux d’événements asynchrones.  
    \item Développer des méthodes de détection d’égomotion adaptées aux données DVS346.  
    \item Intégrer ces algorithmes sur le robot Limo et tester leurs performances dans des scénarios réels de navigation.  
    \item Comparer les performances avec des méthodes basées sur des caméras classiques, en termes de latence, précision et robustesse.  
    \item Explorer l’utilisation de capteurs complémentaires (IMU, lidar) pour améliorer la perception et la navigation.  
\end{itemize}

\section{Structure du rapport}

Le rapport est structuré comme suit :  
\begin{itemize}
    \item \textbf{Chapitre 2 : État de l’art} — présentation des capteurs neuromorphiques, méthodes classiques et récentes pour la navigation robotique.  
    \item \textbf{Chapitre 3 : Matériel et méthodes} — description du robot Limo, de la caméra DVS346, des algorithmes de perception et des outils logiciels utilisés.  
    \item \textbf{Chapitre 4 : Expérimentations et résultats} — protocole expérimental, résultats obtenus, analyse et comparaison.  
    \item \textbf{Chapitre 5 : Conclusion et perspectives} — synthèse des travaux et directions futures pour la vision neuromorphique en robotique mobile.  
\end{itemize}
\end{comment}

%\section*{Introduction}
Les caméras neuromorphiques (ou caméras à événements) sont des capteurs bio-inspirés qui fonctionnent différemment des caméras traditionnelles. Plutôt que de capturer des images à fréquence fixe, elles n'enregistrent que les changements de luminosité au niveau de chaque pixel, de manière asynchrone \cite{gallego2020event}. Chaque « événement » correspond à une variation brusque d'intensité lumineuse en un point, horodatée avec une précision microsecondes et associée à une polarité (augmentation ou diminution de luminosité) \cite{bugueno2025human}. 

Ce mode de fonctionnement événementiel présente plusieurs atouts majeurs : latence ultra-faible (de l'ordre de la microseconde), haute résolution temporelle, très grande plage dynamique (typiquement ~120 dB, contre ~60 dB pour une caméra standard), faible consommation et quasi absence de flou de mouvement, même pour des objets rapides ou en conditions d'éclairage extrêmes \cite{zhu2023overleaf, sanket2020evdodgenet}. 

Le modèle DAVIS346 d'iniVation, objet de cette étude, illustre bien ces propriétés : il combine un capteur à événements (DVS) de résolution 346×260 et un capteur d'intensité APS global shutter, avec une latence des événements de l'ordre de ~20 µs et une dynamique supérieure à 120 dB \cite{gallego2020event}. Il intègre en outre un IMU 6 axes à 1 kHz, ce qui permet la fusion temporelle des événements avec des données inertiales pour des tâches de navigation avancées.

\section{Caméras neuromorphiques et vision événementielle}
\subsection{Principe du capteur à événements (DAVIS346) et avantages}
Une caméra neuromorphique comme la DAVIS346 fonctionne selon le principe de la vision événementielle. Chaque pixel du capteur agit indépendamment et déclenche un événement dès qu'il détecte un changement local de luminosité dépassant un certain seuil. Un événement est généralement défini par : les coordonnées du pixel $(x,y)$, le timestamp $t$ (avec une résolution temporelle pouvant atteindre la microseconde), et une polarité $p$ indiquant si le changement est une augmentation (+1) ou une diminution (-1) d'intensité \cite{bugueno2025human}.

Les avantages de ce fonctionnement neuromorphique sont multiples :
\begin{itemize}
\item \textbf{Latence extrêmement faible} : de l'ordre de quelques microsecondes seulement \cite{zhu2023dynamic}
\item \textbf{Haute résolution temporelle} : fréquence équivalente d'échantillonnage très élevée
\item \textbf{Grande plage dynamique} : typiquement >120 dB contre ~60-70 dB pour une caméra CMOS standard
\item \textbf{Basse consommation et encombrement} : consommation typique de quelques dizaines de milliwatts
\item \textbf{Robustesse au flou de mouvement} : élimination du flou cinétique grâce au déclenchement indépendant des pixels
\end{itemize}

\subsection{Périphérique cible : plateforme Agilex LIMO sous ROS 2 et Gazebo}
Le projet s'articule autour du robot mobile Agilex LIMO, équipé de la caméra DAVIS346. Deux contextes sont envisagés :
\begin{itemize}
\item \textbf{Simulation} : intégration d'un modèle virtuel du capteur événementiel dans Gazebo \cite{hbp2024gazebo}
\item \textbf{Expérimentation réelle} : intégration physique avec driver ROS 2 \cite{rosevent2024libcaer}
\end{itemize}

Pour la simulation, on utilise le plugin DVS Gazebo développé par le Human Brain Project \cite{hbp2024gazebo}. Pour l'intégration réelle, le driver recommandé est \texttt{libcaer\_driver} \cite{rosevent2024libcaer}.

\section{Perception par caméras neuromorphiques pour la navigation autonome}
\subsection{Évitement d'obstacles avec caméras à événements}
\begin{itemize}
\item \textbf{Obstacles statiques} : Approche monoculaire par prédiction de carte de profondeur via réseau de neurones \cite{bhattacharya2024monocular}
\item \textbf{Obstacles dynamiques} : Système d'évitement de balles et objets lancés \cite{falanga2020dynamic} avec double caméra événementielle
\item \textbf{Robotique terrestre} : Combinaison caméra événementielle et RGB-D pour éviter des objets lancés \cite{zhu2023dynamic}
\end{itemize}

\subsection{Estimation d'égomotion et SLAM événementiel}
\begin{itemize}
\item \textbf{Méthodes géométriques} : Maximisation du contraste, extraction de features asynchrones, fusion événements+images+IMU \cite{rebecq2017evo, rosinol2018ultimate}
\item \textbf{Méthodes par apprentissage} : Réseaux de neurones pour flux optique (EV-FlowNet), odométrie 6-DOF, SNN pour traitement asynchrone \cite{gehrig2020eklt}
\end{itemize}

\subsection{Navigation autonome guidée par vision neuromorphique}
\begin{itemize}
\item \textbf{Contrôle réactif bio-inspiré} : Suivi de mur et évitement par régulation du flux optique
\item \textbf{Approches par apprentissage et RL} : Politiques de navigation apprises à partir de représentations événementielles \cite{bugueno2025human}
\item \textbf{SLAM et navigation globale} : Intégration dans des pipelines SLAM complets avec keyframes événementielles \cite{rebecq2017evo}
\end{itemize}

\section{Algorithmes pour les données événementielles : approches classiques vs apprentissage}
\subsection{Représentation et pré-traitements des événements}
\begin{itemize}
\item Images d'accumulation (event frames) \cite{bugueno2025human}
\item Surface d'événements actifs (SAE)
\item Grille spatio-temporelle (voxel grid)
\item Événements reconstruits en images \cite{scheerlinck2020fast}
\end{itemize}

\subsection{Approches classiques (non apprenantes)}
\begin{itemize}
\item Détection de features événementielles (eFAST, Harris3D)
\item Flux optique et suivi d'objets \cite{zhu2023dynamic}
\item Cartographie et localisation classiques (filtres Bayésiens asynchrones, ICP)
\end{itemize}

\subsection{Approches par apprentissage profond (Deep Learning)}
\begin{itemize}
\item \textbf{ANN conventionnels} : Classification, détection d'objets, estimation de mouvement, pilotage end-to-end \cite{delbruck2018endtoend}
\item \textbf{Réseaux de neurones à spikes (SNN)} : Architectures neuromorphiques pour ultra-basse consommation \cite{kaiser2016towards}
\end{itemize}

\section{Exemples de projets et cas d'usage en robotique mobile}
\begin{itemize}
\item \textbf{Drone racer autonome (UZH, 2019)} : Course de drones en environnement intérieur avec DAVIS240 \cite{mueggler2017event}
\item \textbf{Voiturette auto pour faibles latences (Prophesee, 2020)} : Détection de piétons en ~50 ms \cite{nature2024lowlatency}
\item \textbf{Robot mobile d'intérieur (RPG, 2017)} : Navigation en grotte sombre avec éclairage dynamique
\item \textbf{TurtleBot avec DVS (ULiege, 2020)} : Détection de mouvements humains pour interaction sécuritaire \cite{csdn2021ros}
\item \textbf{Projet CERBERE (France, 2022)} : Détection d'objets rapides sur véhicule autonome \cite{anr2022cerbere}
\end{itemize}

\section{Intégration sous ROS 2 et simulation Gazebo : ressources et meilleures pratiques}
\subsection{Drivers et messages ROS 2}
\begin{itemize}
\item Driver \texttt{libcaer\_driver} pour communication avec la DAVIS346 \cite{rosevent2024libcaer}
\item Messages standardisés dans \texttt{event\_camera\_msgs} \cite{rosevent2024eventcamera}
\item Visualisation via conversion en images OpenCV
\end{itemize}

\subsection{Simulation Gazebo (Ignition)}
\begin{itemize}
\item Intégration via plugin DVS dans le URDF du LIMO \cite{hbp2024gazebo}
\item Configuration des paramètres (seuil de contraste, fréquence de rendu)
\item Alternatives : ESIM \cite{mueggler2017event}, Microsoft AirSim
\end{itemize}

\subsection{Ressources et limitations}
\begin{itemize}
\item Packages ROS existants : \texttt{rpg\_dvs\_ros} \cite{rpg2024dvs}, \texttt{esvo\_core} \cite{hong2024esvo}, \texttt{metavision\_ros} \cite{prophesee2024metavision}
\item Limitations actuelles : communauté ROS2 naissante, visualisation limitée dans RViz2
\end{itemize}

\section{Propositions de projet, méthodologie de recherche et plan de travail}
\subsection{Idées de projets concrets réalisables en ~100 heures}
\begin{enumerate}
\item \textbf{Projet A} – Évitement de collisions dynamiques en entrepôt
\item \textbf{Projet B} – Suivi de trajectoire à grande vitesse par flux optique neuromorphique
\item \textbf{Projet C} – Odométrie visuelle-inertielle neuromorphique pour localisation en conditions dégradées
\end{enumerate}

\subsection{Méthodologie de recherche rigoureuse et structurée}
\begin{itemize}
\item Définition précise des objectifs et hypothèses
\item État de l'art ciblé
\item Planification par étapes avec jalons mesurables
\item Approche expérimentale et itérative
\item Documentation continue
\item Validation rigoureuse
\item Analyse critique des résultats
\item Communication et reproduction
\end{itemize}

\subsection{Plan de travail progressif}
\begin{enumerate}
\item \textbf{Phase 1} : Mise en place de la simulation et tests préliminaires (20h)
\item \textbf{Phase 2} : Développement de l'algorithme de perception (40h)
\item \textbf{Phase 3} : Intégration sur la plateforme réelle (30h)
\item \textbf{Phase 4} : Validation expérimentale et documentation finale (10h)
\end{enumerate}

\section*{Conclusion}
Les caméras neuromorphiques telles que la DAVIS346 représentent une technologie prometteuse pour la robotique mobile, en offrant une perception ultra-rapide et robuste là où la vision classique montre ses limites. Cette étude bibliographique a présenté un état de l'art complet des méthodes et applications, démontrant le potentiel de ces capteurs pour améliorer les capacités de navigation des robots mobiles comme le LIMO.

